\begin{afterword}
\summaryhead

The essay argues that people reason badly about politics because, in the environment humans evolved in, political disputes could decide who lived and who died. Political arguments become like soldiers: once you choose a side, you defend every argument for it and attack every argument against it.

For that reason the author advises against using contemporary political examples when teaching or discussing rationality; if a point is about politics, a historical case such as Louis XVI will do. The illustration is a standard problem in artificial intelligence about Nixon, a Quaker and a Republican, which the author thinks needlessly takes a side. The post does not ask for readers to be apolitical, only to avoid gratuitous jabs at a party, because some readers may belong to it and because this is better for the community.

\responsehead

The essay is short, and most of its length is not argument. Its evidence for the central claim is an evolutionary story declared ``obvious'' and then not examined. Its one worked example, the Nixon problem, is misstated so that it no longer shows what nonmonotonic reasoning is about. Two of its eight paragraphs attribute a motive, hedged with ``probably,'' to the example's unnamed author, which is the kind of dig the essay goes on to discourage.

The advice that remains is a rule of etiquette: when teaching, avoid needless partisan examples. It is not what the title says. The title says politics kills minds. The body says only that politics is a poor place to practice. The escape clause, ``unless all the discussants are already rational,'' leaves every group free to decide the rule does not apply to it.

The stated reason for the rule is revealing. ``It doesn't matter whether (you think) the Republican Party really is at fault,'' because avoiding the charge is ``better for the spiritual growth of the community.'' The rule protects a community's cohesion, and it says so. Cohesion is a different aim from finding out what is true. The essay's opening argument, that politics degrades reasoning, is a reason about truth; here the post argues only for cohesion.

The title is what readers met. In 2014 Rob Bensinger collected objections from people near the community. One of them, Miri Mogilevsky, wrote that most outsiders ``initially encounter `politics is the mind-killer' being thrown out in comment threads,'' used to dismiss people who raised political concerns. In the same discussion the author wrote that with ``national politics as seen on TV'' in place of ``politics'' in many places, ``it would have more precisely conveyed what I was trying to say.'' The narrower claim is the one the author meant. The broad one is the one printed, and it is the version in the Highlights collection.

In practice, the essay gave a community a phrase that ends a conversation. It tells members that political disagreement is a symptom, and it lets any member diagnose it in others without addressing their argument. That is an inference, and the 2014 accounts above support it. The essay's own analysis, arguments as soldiers, describes what the phrase became.

In short: an etiquette rule for a forum, justified by an untested evolutionary story and a misstated example, under a title that claims far more than the text and that its readers used to end arguments.
\end{afterword}
